\documentclass[11pt]{article}
\usepackage[margin=0.85in]{geometry}
\usepackage{amsmath,amssymb,amsthm,hyperref,booktabs}
\hypersetup{colorlinks=true,urlcolor=blue,linkcolor=blue}
\newtheorem{theorem}{Theorem}
\setlength{\emergencystretch}{3em}
\title{A certified counterexample to complete monotonicity\\
\large MathDB 351688: the Mills-ratio derivative quotient}
\author{Principia Math In Progress}
\date{17 September 2026; source audit updated 22 September 2026}
\begin{document}
\maketitle

\section*{Task 1. Resolve the stated complete-monotonicity conjecture}
Let
\[
 r(z)=e^{z^2/2}\int_z^\infty e^{-t^2/2}\,dt,
 \qquad F_n(z)=-\frac{(r^{(n)}(z))^2}{r^{(2n+1)}(z)}.
\]
Conjecture 3 in Ferreira--Simon \cite[p.~31]{FS} asserts that every
\(F_n\) is completely monotone on \((0,\infty)\). The definition of
\(r\) is on p.~29. Complete monotonicity requires
\((-1)^kF_n^{(k)}(z)\geq0\) for every nonnegative integer \(k\) and
\(z>0\). The following counterexample uses \(n=1\).

\begin{theorem}
For the exact positive argument \(c=1/100\),
\begin{equation}\label{negative}
 -\frac{461}{10^{54}}
 < \frac{F_1^{(70)}(c)}{70!}
 < -\frac{460}{10^{54}}<0.
\end{equation}
Consequently, \(F_1\) is not completely monotone, and the universal
assertion in Conjecture 3 is false.
\end{theorem}

\noindent\textbf{Publication scope.}
The complete 33-page author manuscript \cite{FS} has now been read.
The authors announced this version as accepted for publication in
Transactions of the American Mathematical Society. The source audit
uses this accepted manuscript; the separately typeset journal version
remains unread. The package records the bounded claim audit.

\begin{proof}
Changing variables in the defining integral yields
\[
 r(z)=\int_0^\infty e^{-zt-t^2/2}\,dt,
 \qquad r'(z)=zr(z)-1.
\]
Differentiation under the integral is valid locally uniformly for real
\(z\), because the Gaussian dominates every polynomial times a fixed
exponential. In particular,
\[
 -r'''(z)=\int_0^\infty t^3e^{-zt-t^2/2}\,dt>0.
\]
Thus the quotient and all its derivatives are well defined on the
required interval.

Put \(R(w)=r(c-w)=\sum_{k\geq0}e_kw^k\). The differential equation
for \(r\) gives the exact coefficient recurrence
\begin{equation}\label{rrec}
 e_0=r(c),\qquad e_1=1-ce_0,\qquad
 e_{k+1}=\frac{e_{k-1}-ce_k}{k+1}\quad(k\geq1).
\end{equation}
Moreover,
\[
 F_1(c-w)=\frac{R'(w)^2}{R'''(w)}=\sum_{k\geq0}q_kw^k,
 \qquad q_k=\frac{(-1)^kF_1^{(k)}(c)}{k!}.
\]
Define
\[
 a_k=(k+1)e_{k+1},\qquad
 b_k=(k+1)(k+2)(k+3)e_{k+3}.
\]
Since \(b_0=-r'''(c)>0\), coefficient comparison gives
\begin{equation}\label{qrec}
 q_k=\frac{1}{b_0}\left(
 \sum_{j=0}^ka_ja_{k-j}-\sum_{j=1}^kb_jq_{k-j}\right).
\end{equation}
Computing through \(e_{73}\) and \(q_{70}\) therefore determines the
required derivative without numerical differentiation.

For completeness, we specify the exact finite certificate for its sign.
Every bound below uses integers and rational numbers. Set
\[
 A(v)=\sum_{j=0}^{199}\frac{(-1)^jv^{2j+1}}{2j+1},
 \qquad T(v)=\frac{v^{401}}{401}.
\]
The alternating-series estimate gives
\(A(v)<\arctan v<A(v)+T(v)\) for \(v=1/5,1/239\).
Machin's identity \(\pi=16\arctan(1/5)-4\arctan(1/239)\) then gives
explicit rational lower and upper bounds \(\pi_-<\pi<\pi_+\).
The identity follows directly from the tangent addition formula;
the angle \(4\arctan(1/5)-\arctan(1/239)\) lies in \((0,\pi/2)\)
and has tangent one.

Let \(S=10^{200}\) and
\[
 h=\left\lfloor S\sqrt{\pi_-/2}\right\rfloor.
\]
Integer comparisons verify
\[
 (h/S)^2<\pi_-/2<\pi_+/2<((h+1)/S)^2.
\]
Consequently \(h/S<\sqrt{\pi/2}<(h+1)/S\).
To enclose \(e_0=r(c)\), use
\[
 r(c)=e^{c^2/2}\left(\sqrt{\pi/2}-\int_0^c e^{-t^2/2}\,dt\right).
\]
For \(u=c^2/2=1/20000\), define
\[
 E=\sum_{j=0}^{79}\frac{u^j}{j!},\qquad
 J=\sum_{j=0}^{79}\frac{(-1)^jc^{2j+1}}{2^jj!(2j+1)}.
\]
Positive-series and alternating-series bounds give
\begin{align*}
 E&<e^u<E+\frac{u^{80}}{80!}\frac1{1-u/81},\\
 J&<\int_0^ce^{-t^2/2}\,dt
   <J+\frac{c^{161}}{2^{80}80!\,161}.
\end{align*}
These enclose the sole nonrational input \(e_0\).

Propagate these intervals through \eqref{rrec} and \eqref{qrec}, rounding
outwards after every operation to the grid \(S^{-1}\mathbb Z\).
For multiplication take the minimum and maximum of the four endpoint
products; for division take the four endpoint quotients, having checked
that the denominator interval is strictly positive. Round lower endpoints
down and upper endpoints up. Induction on the recurrences proves that
every computed interval contains the exact coefficient.
The accompanying standard-library Python verifier performs precisely
these operations, using arbitrary-length integers. Its final interval
is strictly inside
\[
 \left(-\frac{461}{10^{54}},-\frac{460}{10^{54}}\right).
\]
This is \eqref{negative}, because the derivative order 70 is even.
The exact rational endpoints are saved in \path{exact-certificate.json}.
This finite certificate proves a counterexample to a universal statement;
it is not a finite-sample argument for an infinite assertion.
\end{proof}

The same verifier additionally certifies
\[
 -\frac{302}{10^{28}}<\frac{F_0^{(40)}(0)}{40!}
 <-\frac{301}{10^{28}},\qquad
 -\frac{758}{10^{54}}<\frac{F_1^{(70)}(0)}{70!}
 <-\frac{757}{10^{54}}.
\]
Continuity would already give violations at small positive arguments.
The theorem instead supplies the explicit point \(1/100\).
This result does not contradict Proposition 10 of \cite{FS}, which
concerns the different functions \(z\mapsto F_n(\sqrt z)\).

\section*{An analytic explanation of the failure}
The failure also has a proof that does not depend on the coefficient
computation. In fact, at every fixed positive argument there are
infinitely many derivative orders with the wrong sign.

Write \(h(z)=-r'(z)=1-zr(z)\) and
\[
 H(z)=-r'''(z)=\int_0^\infty t^3e^{-zt-t^2/2}\,dt
              =(z^2+3)h(z)-1.
\]
Both functions are entire, and \(H(x)>0\) for every real \(x\).
Moreover,
\[
 |H(z)|\leq8e^{|z|^2},\qquad
 x^4H(x)=\int_0^\infty s^3e^{-s-s^2/(2x^2)}\,ds\longrightarrow6
 \quad(x\to+\infty).
\]
The growth bound follows from
\(|z|t\leq |z|^2+t^2/4\); the limit follows by dominated convergence.

Suppose \(H\) had no zeros. It would have an entire logarithm \(g\),
with \(H=e^g\) and \(\Re g(z)\leq\log8+|z|^2\).
For the Taylor coefficients \(g(z)=\sum d_jz^j\), the Fourier
coefficient formula on the circle of radius \(R\), applied to
\(\log8+R^2-\Re g\geq0\), gives
\[
 |d_j|R^j\leq2(\log8+R^2-\Re g(0)),\qquad j\geq1.
\]
Letting \(R\to\infty\) shows that \(g\) has degree at most two.
But \(\log|H(x)|=\Re g(x)\) would then be a real polynomial of
degree at most two, whereas the displayed limit gives
\(\log H(x)=-4\log x+\log6+o(1)\). This is impossible.
Hence \(H\) has a nonreal zero. At every zero of \(H\), the identity
\(H=(z^2+3)h-1\) implies \(h\ne0\). Thus \(F_1=h^2/H\) has a
nonreal pole and has no real poles.

We recall a short form of the positive-coefficient argument.
If \(G(w)=\sum a_jw^j\), \(a_j\geq0\), has finite positive
radius of convergence \(\rho\), it cannot be analytic at \(w=\rho\).
Indeed, if it were, monotone convergence as \(w\uparrow\rho\) would give
\[
 \frac{G^{(k)}(\rho)}{k!}
 =\sum_{j\geq k}\binom jk a_j\rho^{j-k}.
\]
The Taylor series at \(\rho\) converges at \(\rho+\varepsilon\)
for some \(\varepsilon>0\). Interchanging its nonnegative sums then
gives \(\sum_j a_j(\rho+\varepsilon)^j<\infty\), a contradiction.

For any fixed \(c>0\), the Taylor series of \(F_1(c-w)\) has finite
positive radius because of its nonreal pole, but the function is
analytic at every real \(w\). Its coefficients therefore cannot all
be nonnegative. Nor can they be eventually nonnegative: subtracting
finitely many Taylor terms would give the same contradiction.
Since these coefficients are \((-1)^kF_1^{(k)}(c)/k!\), infinitely
many have the wrong sign. This proves non-complete-monotonicity
analytically, as well as explaining why a high derivative can detect it.

\section*{Tasks 2--5. Proposed directions after the counterexample}
These are proposed mathematical tasks, not additional solved claims.
\begin{description}
\item[Task 2.] Determine the smallest order \(k\) for which
\((-1)^kF_n^{(k)}(z)<0\) somewhere on \((0,\infty)\), first for
\(n=0,1\), and then its dependence on \(n\).
\item[Task 3.] Decide whether every \(F_n\), \(n\geq0\), fails complete
monotonicity. Extend the question to real moment orders by replacing
\((-1)^nr^{(n)}(z)\) with
\(\int_0^\infty t^\nu e^{-zt-t^2/2}\,dt\), \(\nu>-1\).
\item[Task 4.] Characterize the powers \(\beta>0\) for which
\(z\mapsto F_n(z^\beta)\) is completely monotone. In particular,
determine whether the known exponent \(\beta=1/2\) is sharp.
\item[Task 5.] Locate the nearest complex zeros of \(r^{(2n+1)}\),
identify cancellations in \(F_n\), and use the resulting singularities
to derive asymptotics and sign changes of its high derivatives.
\end{description}

\begin{thebibliography}{9}
\bibitem{FS} R. A. C. Ferreira and T. Simon,
\emph{Convolution of beta prime distribution},
arXiv:2108.09244v2 (19 May 2022), Section 5.5, Conjecture 3,
\url{https://arxiv.org/abs/2108.09244v2}.
Journal metadata: Trans. Amer. Math. Soc. 376 (2023), 855--890,
\url{https://doi.org/10.1090/tran/8748}.
The author-announced accepted manuscript was read completely;
the separately typeset version of record was not accessible.
\end{thebibliography}
\end{document}
